\section{Opening referrals after provisional offers}
Before allocation a recommendation changes exposure. After provisional offers an intermediary may have to reconcile full-market stability with promised ranks. Suppose each participant has a fixed representative and a paid audit opens every feasible edge incident to that representative's clients. Let $M_0$ be DA on the initial visibility graph and $M^+$ applicant-proposing DA under full visibility. Each applicant who prefers $M_0$ to $M^+$ must lose rank under every full-stable repair. For a protected group, a planner can spare every other member simultaneously by full disclosure. The \emph{protection premium} is the least cost of obtaining this guarantee with a globally stable DA outcome, minus the least cost of global stability alone. These actions and costs differ from the one-slot display edits studied before offers.

Theorems~\ref{thm:two} and~\ref{thm:three} give a sharp one-blocker boundary. If everyone has at most two acceptable partners and the provisional matching has at most one full-market blocker, every successful repair already protects the group. At degree three, for every odd $n=2m+1$ there is a connected market with one client per firm, unit costs, a perfect provisional assignment and exactly one blocker in which one audit repairs stability but protecting all applicants costs $m+1$. The mechanism is a rejection chain: opening one hidden offer gives the receiving side its favorite partners and moves all but one applicant below their original ranks. Protecting those ranks requires opening $m+1$ disjoint hidden opportunities. The appendix proves both results, gives the linear-time private-firm algorithm at degree two, and establishes hardness when a representative can serve multiple clients.

In a separately specified audit benchmark with independent listing lotteries, 40 blocks and four boroughs, lower-income protection cost nothing extra in all 480 markets under uniform top-five lists and under cheapest-first lists. Its donor split, catalogue cleaning and utility specification differ from the calibrated simulations. In one market, 15 audits made all edges of a full-stable allocation visible but all 4,096 minimum covers returned the same blocked DA outcome. Both ordinary and protected repair required 16 audits. A zero premium means a cheapest plan exists that meets the guarantee. It says little about how many applicants had an offer to protect.
